\chapter{Plataforma web en la nube}

\section{Propósito del capítulo}

Este capítulo documenta la plataforma web que constituye la capa de interfaz
humana y de coordinación central del sistema de monitoreo espectral ANE. La
plataforma está compuesta por un frontend de aplicación de página única (SPA)
y un backend de servicios que actúa como centro de orquestación entre las
personas operadoras, los sensores RF en campo, la base de datos y los
servicios de postprocesamiento.

El capítulo describe la arquitectura general, las funcionalidades ofrecidas a
los usuarios, las interfaces de comunicación entre componentes, las decisiones
de diseño que guiaron la implementación, los problemas encontrados durante el
desarrollo, los aprendizajes obtenidos y las limitaciones actuales. El
objetivo es proporcionar una comprensión completa de cómo la plataforma web
hace posible el ciclo operativo completo: desde la autenticación de una
persona usuaria hasta la generación de un reporte de cumplimiento normativo
basado en mediciones de espectro capturadas en campo.

\section{Contexto dentro del sistema general}

La plataforma web ocupa la posición central en la arquitectura del sistema
ANE. Como se describe en el capítulo de diseño general (Capítulo~3), el
sistema se organiza en tres dominios principales: borde (sensores RF y
adquisición), nube (plataforma web, base de datos y postprocesamiento) y
usuarios (personas operadoras y administradoras). La plataforma web constituye
la totalidad del dominio de nube desde la perspectiva de coordinación y la
totalidad del dominio de usuarios desde la perspectiva de interfaz.

La relación con los demás módulos del sistema es la siguiente:

\begin{itemize}
    \item \textbf{Hardware y sensores RF (Capítulo~4)}: la plataforma recibe
    los datos de espectro, estado, GPS y audio demodulado que los sensores
    capturan en campo. También entrega a los sensores las configuraciones de
    adquisición y las campañas programadas que deben ejecutar.

    \item \textbf{SDR y adquisición en borde (Capítulo~5)}: el flujo de
    adquisición definido en el borde depende de los parámetros que la
    plataforma configura y transmite. La plataforma es la consumidora de los
    datos generados por la cadena de procesamiento SDR.

    \item \textbf{Análisis espectral y localización en la nube (Capítulo~7)}:
    la plataforma invoca al servicio de postprocesamiento Python para el
    análisis de emisiones, detección de señales, cruce con licencias y
    evaluación de cumplimiento normativo. Los resultados del análisis se
    integran en los reportes que la plataforma entrega a las personas usuarias.

    \item \textbf{Protocolos de prueba (Capítulo~8)}: la plataforma es el
    objeto de pruebas de integración extremo a extremo y pruebas de usabilidad.
    Los protocolos de prueba definidos en ese capítulo se aplican sobre la
    plataforma para verificar su correcto funcionamiento en conjunto con los
    demás módulos.

    \item \textbf{Integración general (Capítulo~9)}: la plataforma es el
    componente integrador por excelencia. El capítulo de integración detalla
    cómo todos los módulos se conectan a través de la plataforma para formar
    el sistema completo.
\end{itemize}

En el flujo extremo a extremo del sistema, la plataforma web es el punto de
entrada para las personas y el punto de convergencia para los datos. Sin la
plataforma, los sensores podrían capturar datos pero no habría forma de
configurarlos, visualizar sus mediciones, planificar campañas, detectar
alertas ni generar reportes.

\section{Desarrollo técnico}

La plataforma web se implementa como dos aplicaciones independientes que se
despliegan como contenedores Docker coordinados por un proxy inverso Nginx:
un frontend de aplicación de página única (SPA) y un backend de servicios API.

\subsection{Frontend}

El frontend es una aplicación web de página única desarrollada con React 19 y
TypeScript, construida con el empaquetador Vite. La interfaz sigue un diseño
de panel de control con menú lateral de navegación que da acceso a siete
pantallas principales.

\subsubsection*{Tecnologías}

La pila tecnológica del frontend está conformada por:

\begin{itemize}
    \item \textbf{React 19} como librería principal de componentes de interfaz,
    seleccionado por su modelo eficiente de actualización del DOM virtual y su
    capacidad para manejar datos en tiempo real sin re-renderizaciones
    innecesarias.

    \item \textbf{TypeScript} como lenguaje de desarrollo, proporcionando
    tipado estático que mejora la detección temprana de errores y documenta
    implícitamente las interfaces de datos entre componentes.

    \item \textbf{Vite 5} como empaquetador y servidor de desarrollo,
    seleccionado por su velocidad de compilación basada en ESM nativo y su
    configuración simplificada.

    \item \textbf{Tailwind CSS 3} como framework de estilos mediante
    utilidades atómicas, facilitando la iteración rápida sobre el diseño y la
    consistencia visual entre componentes.

    \item \textbf{Plotly.js} para la visualización interactiva de espectro de
    potencia y gráficas waterfall (espectrograma), seleccionada por su soporte
    nativo para gráficos científicos con miles de puntos de datos y
    actualización en tiempo real.

    \item \textbf{Leaflet} y \textbf{React-Leaflet} para la visualización de
    mapas geográficos con marcadores de sensores, permitiendo la ubicación
    espacial de la red de monitoreo.

    \item \textbf{Recharts} para gráficas estadísticas complementarias en el
    panel de inicio.

    \item \textbf{React Router 7} para el enrutamiento del lado del cliente,
    permitiendo navegación entre pantallas sin recarga completa de la página.

    \item \textbf{Microsoft Authentication Library (MSAL)} para la integración
    con Azure AD como proveedor de identidad institucional.

    \item \textbf{Axios} como cliente HTTP para la comunicación con el backend.
\end{itemize}

\subsubsection*{Pantallas principales}

La interfaz se organiza en las siguientes pantallas, accesibles desde un menú
lateral cuya visibilidad se adapta al rol de la persona usuaria:

\textbf{Inicio (Dashboard)}: pantalla de bienvenida que presenta un mapa
geográfico con la ubicación de todos los sensores registrados, indicadores de
estado por colores, y un resumen de estadísticas del sistema (sensores
activos, campañas en curso, alertas recientes). Es el punto de partida para la
navegación.

\textbf{Dispositivos}: vista de gestión de la red de sensores. Muestra el
listado completo de sensores con su estado actual, ubicación y antenas
asociadas. Incluye un mapa interactivo con marcadores de estado y acceso al
detalle de cada sensor. Desde esta pantalla se puede iniciar un monitoreo en
vivo sobre un sensor seleccionado.

\textbf{Monitoreo en vivo}: flujo central de la plataforma para realizar
mediciones inmediatas. La persona selecciona un sensor y una antena, define
los parámetros de adquisición (rango de frecuencia, ancho de banda,
resolución, ganancias) e inicia la captura. Durante la medición se visualiza
en tiempo real el espectro de potencia (PSD) mediante Plotly.js, la evolución
temporal en waterfall, y el audio demodulado AM/FM cuando está habilitado. La
persona puede detener la medición o guardar la configuración como una nueva
campaña.

\textbf{Campañas}: gestión de mediciones programadas. Permite crear campañas
definiendo nombre, fechas de inicio y fin, horas de operación, intervalo entre
mediciones, parámetros espectrales y sensores asignados. Incluye listado con
filtros por estado, visualización de detalle, inicio y detención manual, y
consulta de datos capturados. Los datos de campaña son la fuente para la
generación de reportes de cumplimiento.

\textbf{Alertas}: revisión de incidentes y eventos operativos detectados por
el sistema. Presenta un listado cronológico con filtros por sensor, tipo de
alerta y rango de fechas. Cada alerta se vincula con los datos del sensor
para facilitar el diagnóstico.

\textbf{Configuración}: panel administrativo para la gestión de recursos.
Incluye CRUD del catálogo de antenas (tipo, rango de frecuencia, ganancia,
código de inventario), administración de cuentas de usuario (roles y
permisos), y configuración de parámetros globales del sistema.

\textbf{Audio}: reproductor de audio demodulado AM/FM, accesible desde el
monitoreo en vivo o desde la vista de un sensor. Muestra métricas de
modulación en tiempo real: profundidad para AM, excursión para FM.

\subsubsection*{Mecanismo de actualización en tiempo real}

El frontend mantiene una conexión WebSocket con el backend para recibir
actualizaciones sin necesidad de recarga manual. Los eventos recibidos
incluyen cambios de estado de sensores, nuevas ubicaciones GPS, datos de
espectro en tiempo real y confirmaciones de comandos. Estos eventos actualizan
automáticamente los mapas, gráficas e indicadores en todas las pantallas
relevantes, implementando un modelo de datos reactivo donde la interfaz
refleja el estado del sistema en todo momento.

\subsection{Backend}

El backend es el centro de coordinación de la plataforma, implementado como un
servidor Node.js con TypeScript sobre el framework Express. Su función
principal es conectar a las personas usuarias (a través del frontend), los
sensores RF en campo, la base de datos PostgreSQL y los servicios de
postprocesamiento.

\subsubsection*{Tecnologías}

\begin{itemize}
    \item \textbf{Node.js 18+} como entorno de ejecución, seleccionado por su
    modelo de E/S no bloqueante que permite manejar múltiples conexiones
    concurrentes (frontend, sensores, WebSocket) en un solo proceso.

    \item \textbf{TypeScript} como lenguaje, compilado a JavaScript para
    producción, proporcionando verificación de tipos en tiempo de compilación.

    \item \textbf{Express} como framework HTTP, ofreciendo un modelo maduro de
    middleware, enrutamiento y manejo de solicitudes.

    \item \textbf{WebSocket nativo} (librería ws) para comunicación en tiempo
    real con el frontend, seleccionado sobre alternativas como Socket.IO por su
    menor sobrecarga de protocolo.

    \item \textbf{PostgreSQL 13+} con cliente pg y pool de conexiones, como
    sistema de base de datos relacional.

    \item \textbf{TimescaleDB} como extensión opcional de PostgreSQL para
    optimización de datos de series temporales.

    \item \textbf{JSON Web Tokens (JWT)} y librería jose para la gestión de
    sesiones y autenticación.

    \item \textbf{Swagger/OpenAPI} para la documentación interactiva de la API.
\end{itemize}

\subsubsection*{Capas de la arquitectura}

El backend está organizado en seis capas lógicas con responsabilidades bien
definidas:

\textbf{Capa de entrada HTTP}: expone el servidor en el puerto configurado
(3000) y aplica middleware común: CORS para permitir solicitudes del frontend,
registro (logging) de cada solicitud con método y ruta, y parsing de JSON con
un límite ampliado de 50 MB para permitir la recepción de tramas espectrales
de alta resolución enviadas por los sensores.

\textbf{Capa de autenticación y autorización}: valida la identidad de la
persona usuaria mediante dos mecanismos complementarios. Para el acceso
institucional desde el frontend, se integra con Azure AD utilizando el
protocolo OpenID Connect y la librería jose para validación de tokens. Para el
acceso administrativo local, se utilizan credenciales almacenadas en la base
de datos con hash bcrypt y tokens JWT firmados por el servidor. Los middleware
de autenticación se aplican de forma selectiva según la ruta: los endpoints de
ingesta de datos desde sensores operan con validación por dirección MAC; los
endpoints de administración requieren sesión validada con rol apropiado.

\textbf{Capa de rutas y controladores}: organiza la lógica de negocio en
módulos independientes: autenticación, gestión de sensores y antenas (CRUD,
asignación), recepción de datos de campo (estado, GPS, espectro, audio),
control remoto de sensores, campañas de medición, reportes de cumplimiento y
configuración del sistema.

\textbf{Capa de modelos y acceso a datos}: abstrae las consultas SQL mediante
funciones tipadas que utilizan un pool de conexiones pg.Pool con hasta 20
conexiones simultáneas. Las operaciones que modifican múltiples tablas se
ejecutan dentro de transacciones PostgreSQL para garantizar consistencia.

\textbf{Capa de comunicación en tiempo real}: mantiene un servidor WebSocket
que transmite cinco tipos de eventos al frontend: \texttt{sensor\_status}
(estados de sensores), \texttt{sensor\_gps} (ubicaciones),
\texttt{sensor\_data} (espectro en tiempo real), \texttt{sensor\_configure} y
\texttt{sensor\_stop} (comandos). Un segundo servidor WebSocket independiente
gestiona el streaming de audio demodulado en formato Opus.

\textbf{Capa de mantenimiento automático}: ejecuta tareas programadas en
intervalos regulares de 30 segundos. La tarea principal valida la antigüedad
del último reporte de cada sensor activo y marca como offline aquellos que
exceden el umbral de inactividad. Los cambios de estado se notifican al
frontend mediante broadcast WebSocket.

\subsubsection*{API REST}

La interfaz de programación del backend expone aproximadamente treinta
endpoints organizados en ocho categorías funcionales, todas documentadas
mediante especificación OpenAPI y accesibles a través de Swagger UI:

\begin{enumerate}
    \item \textbf{Sensor Data}: endpoints POST para que los sensores envíen
    estado operativo, ubicación GPS, datos de espectro radioeléctrico y audio
    demodulado. Son los endpoints de mayor tráfico durante la operación normal.

    \item \textbf{Sensor Query}: endpoints GET para consulta de datos
    históricos: último estado, última ubicación, últimos datos de espectro,
    datos por rango temporal y configuración activa, identificando cada sensor
    por su dirección MAC.

    \item \textbf{Sensor Control}: endpoints POST para enviar comandos a los
    sensores: configuración de parámetros de escaneo, comando de detención y
    almacenamiento de configuración activa.

    \item \textbf{Sensors Management}: CRUD completo de sensores más consultas
    por ID o MAC, y gestión de la relación sensor-antena (asignar, desasignar,
    consultar).

    \item \textbf{Antennas Management}: CRUD completo del catálogo de antenas.

    \item \textbf{Campaigns}: gestión de campañas con operaciones de creación,
    consulta, actualización, eliminación, inicio, detención y consulta de
    datos capturados.

    \item \textbf{Reports}: generación de reportes de cumplimiento normativo
    que integran datos de campaña, ubicación y análisis del servicio de
    postprocesamiento.

    \item \textbf{System}: endpoint raíz con versión de la API y directorio de
    endpoints.
\end{enumerate}

\subsection{Base de datos}

La plataforma utiliza PostgreSQL como sistema de gestión de base de datos
relacional. El esquema está organizado en cuatro dominios funcionales que
agrupan quince tablas:

\textbf{Dominio de identidad y acceso}: almacena credenciales y roles de las
personas usuarias (tabla \texttt{users}, con hash bcrypt para contraseñas) y
un registro de auditoría de acciones sensibles (tabla \texttt{audit\_logs}).

\textbf{Dominio de sensores y configuración}: mantiene el registro maestro de
sensores RF identificados por MAC única (tabla \texttt{sensors}), el catálogo
de antenas con sus parámetros técnicos (\texttt{antennas}), la relación
muchos-a-muchos entre sensores y antenas con puerto de conexión
(\texttt{sensor\_antennas}), y los parámetros de adquisición por sensor:
frecuencias, sample rate, ganancias, ventana, solapamiento y configuración de
demodulación (\texttt{sensor\_configurations}).

\textbf{Dominio de datos de campo}: almacena las series temporales generadas
por los sensores: estado operativo con métricas de CPU, RAM, temperatura y
latencia (\texttt{sensor\_status}), ubicaciones GPS (\texttt{sensor\_gps}),
datos de espectro con arreglos de potencia, métricas de demodulación AM/FM y
asociación opcional a campaña (\texttt{sensor\_data}), y alertas operativas
históricas (\texttt{sensor\_history\_alert}).

\textbf{Dominio de campañas y reportes}: define las campañas de medición con
parámetros de programación y configuración espectral (\texttt{campaigns}), la
asignación de sensores a campañas (\texttt{campaign\_sensors}), el caché de
reportes de cumplimiento en formato JSON (\texttt{compliance\_reports\_cache}),
y parámetros globales del sistema en formato clave-valor
(\texttt{system\_configurations}).

La tabla \texttt{sensor\_status}, que recibe actualizaciones periódicas de
cada sensor y es la de mayor crecimiento, está configurada como hypertable de
TimescaleDB, particionada por la columna \texttt{timestamp\_ms}. Esta
configuración permite compresión y retención automática de datos, y optimiza
consultas temporales mediante índices compuestos por MAC y timestamp.

\subsection{Despliegue}

La plataforma se despliega en un servidor institucional privado de la ANE,
accesible a través de la red interna de la entidad. La arquitectura de
despliegue utiliza contenedores Docker coordinados por un proxy inverso Nginx:

\begin{itemize}
    \item \textbf{Frontend}: imagen Docker construida en dos etapas:
    compilación con Node.js Alpine y servidor web con Nginx Alpine. Expone el
    puerto 80 con configuración para aplicaciones de página única (SPA) y
    compresión gzip.

    \item \textbf{Backend}: imagen Docker construida en dos etapas:
    compilación de TypeScript y ejecución con Node.js Alpine. Expone el puerto
    3000.

    \item \textbf{Nginx}: actúa como punto único de entrada, enrutando el
    tráfico HTTP al frontend (archivos estáticos) y al backend (API REST en
    \texttt{/api/}, WebSocket en \texttt{/ws/}). La configuración incluye
    timeouts extendidos de una hora para la generación de reportes de larga
    duración.

    \item \textbf{PostgreSQL}: base de datos con extensión TimescaleDB,
    ejecutándose en el mismo servidor. En producción se recomienda la imagen
    \texttt{timescale/timescaledb:latest-pg14}.
\end{itemize}

La configuración de entorno se gestiona mediante variables de entorno
(\texttt{DB\_HOST}, \texttt{DB\_PORT}, \texttt{DB\_NAME}, \texttt{DB\_USER},
\texttt{DB\_PASSWORD}, \texttt{PORT}) que permiten desplegar la misma imagen
en entornos de desarrollo, pruebas y producción sin recompilación.

\section{Entradas, salidas e interfaces}

\subsection{Entradas}

La plataforma recibe datos de tres fuentes principales:

\textbf{Desde el frontend}: solicitudes HTTP de las personas usuarias para
autenticarse, consultar información, configurar mediciones, gestionar
campañas, revisar alertas y generar reportes. Estas solicitudes incluyen
tokens de autenticación (JWT o Azure AD) y parámetros específicos según la
operación.

\textbf{Desde los sensores RF}: datos de campo en formato JSON enviados
mediante solicitudes POST:
\begin{itemize}
    \item Estado del sensor: métricas de CPU por núcleo, RAM, disco,
    temperatura, latencia de red y logs del sistema operativo del dispositivo.
    \item Ubicación GPS: latitud, longitud y altitud.
    \item Datos de espectro: arreglo de potencias PSD (\texttt{Pxx}),
    frecuencia inicial y final en Hz, marca de tiempo en epoch milliseconds,
    coordenadas de captura, y métricas de demodulación AM (profundidad) o FM
    (excursión).
    \item Audio demodulado: tramas de audio PCM codificadas en base64 o en
    formato Opus.
\end{itemize}

\textbf{Desde el servicio de postprocesamiento}: resultados de análisis de
emisiones en formato JSON, con frecuencias centrales detectadas, anchos de
banda, potencias, cruce con licencias y evaluación de cumplimiento por
municipio (código DANE).

\subsection{Salidas}

La plataforma entrega:

\textbf{Al frontend}: respuestas JSON a todas las consultas y operaciones,
eventos WebSocket para actualización en tiempo real de estado de sensores,
GPS, datos de espectro y audio.

\textbf{A los sensores RF}: configuraciones de adquisición (frecuencias,
sample rate, ganancias, ventana, solapamiento, demodulación) y listas de
campañas asignadas con todos los parámetros necesarios para ejecutar
mediciones programadas.

\textbf{A las personas usuarias}: reportes de cumplimiento normativo
estructurados con emisiones detectadas, comparación contra licencias,
evaluación por parámetro (frecuencia central, ancho de banda, potencia) y
estimación de intensidad de campo.

\subsection{Interfaces}

Las interfaces de comunicación entre componentes son:

\begin{itemize}
    \item \textbf{Frontend -- Backend}: API REST sobre HTTP con formato JSON.
    Autenticación mediante tokens JWT en cabecera \texttt{Authorization}.
    WebSocket para eventos en tiempo real y streaming de audio.

    \item \textbf{Sensores -- Backend}: API REST sobre HTTP con formato JSON.
    Identificación del sensor por dirección MAC. Los sensores actúan como
    clientes HTTP que envían datos (POST) y consultan configuración (GET).

    \item \textbf{Backend -- PostgreSQL}: conexión TCP nativa mediante el
    cliente pg sobre el puerto 5432, utilizando un pool de conexiones
    configurable.

    \item \textbf{Backend -- Postprocesamiento}: invocación HTTP síncrona al
    servicio Python Flask expuesto en el puerto 8000, con formato JSON para
    solicitudes y respuestas.

    \item \textbf{Nginx -- Frontend}: entrega de archivos estáticos (HTML, JS,
    CSS, assets) por HTTP en puerto 80.

    \item \textbf{Nginx -- Backend}: proxy inverso HTTP y WebSocket, con
    timeouts extendidos para operaciones de larga duración.
\end{itemize}

\section{Decisiones de diseño o implementación}

\subsection{Separación frontend/backend con API REST}

Se decidió separar el frontend y el backend como aplicaciones independientes
comunicadas por API REST, en lugar de una aplicación monolítica con renderizado
en el servidor. Esta decisión se fundamenta en tres consideraciones:

\begin{enumerate}
    \item \textbf{Despliegue independiente}: el frontend (archivos estáticos) y
    el backend (servidor Node.js) pueden escalarse, actualizarse y
    desplegarse de forma independiente. El frontend se sirve desde Nginx como
    archivos estáticos con caché agresiva; el backend maneja la lógica de
    negocio y la concurrencia.

    \item \textbf{Separación de responsabilidades}: el equipo de frontend
    puede iterar sobre la interfaz sin afectar la lógica del backend, y
    viceversa. La API actúa como contrato estable entre ambas capas.

    \item \textbf{Cliente único vs. múltiples clientes}: aunque actualmente el
    frontend web es el único cliente, la arquitectura REST permite que en el
    futuro otros clientes (aplicaciones móviles, scripts de automatización,
    dashboards externos) consuman la misma API sin modificar el backend.
\end{enumerate}

\subsection{WebSocket sobre polling para tiempo real}

La elección de WebSocket como mecanismo de comunicación en tiempo real, en
lugar de polling HTTP periódico, se justifica por el perfil de tráfico del
sistema:

\begin{itemize}
    \item Durante un monitoreo activo, un sensor puede enviar datos de
    espectro cada pocos segundos. Con polling, el frontend tendría que
    consultar al backend repetidamente para detectar nuevos datos, generando
    tráfico innecesario y latencia.

    \item WebSocket permite que el backend notifique proactivamente al
    frontend solo cuando hay datos nuevos, en un modelo push eficiente.

    \item La librería ws (WebSocket nativo) fue seleccionada sobre Socket.IO
    por su menor sobrecarga de protocolo: no requiere el handshake adicional
    de Socket.IO ni sus mecanismos de fallback a long-polling, que son
    innecesarios en un entorno de red controlada.
\end{itemize}

\subsection{PostgreSQL y TimescaleDB para datos de series temporales}

La decisión de adoptar PostgreSQL como base de datos principal —y
específicamente TimescaleDB como extensión para tablas de series temporales—
fue la de mayor impacto en la arquitectura del backend. El sistema comenzó
utilizando SQLite, una base de datos sin servidor que almacena toda la
información en un solo archivo.

SQLite ofrecía ventajas iniciales de simplicidad: sin necesidad de instalar
ni configurar un servidor de base de datos separado, sin autenticación, con
despliegue trivial. Sin embargo, a medida que el volumen de datos creció
—particularmente en las tablas de estado de sensores y datos de espectro— se
manifestaron limitaciones críticas:

\begin{itemize}
    \item El límite práctico de aproximadamente 2 GB por archivo de base de
    datos comenzó a causar degradación y bloqueos.
    \item SQLite no soporta concurrencia real de escritura: múltiples
    sensores enviando datos simultáneamente generaban contención de bloqueo.
    \item Las consultas analíticas sobre grandes volúmenes de datos
    históricos (necesarias para reportes) tenían rendimiento deficiente.
\end{itemize}

La migración a PostgreSQL resolvió estas limitaciones y habilitó capacidades
nuevas: tipos de datos nativos para valores geoespaciales y series temporales,
transacciones ACID completas en operaciones multi-tabla, índices compuestos
para consultas por sensor y rango temporal, y compresión y retención
automática de datos antiguos mediante TimescaleDB. La migración se realizó de
forma progresiva, actualizando primero la capa de conexión, luego los modelos
de datos y finalmente las rutas con consultas directas.

\subsection{Microservicio de postprocesamiento en Python}

El análisis espectral avanzado —detección de emisiones, estimación de piso de
ruido por escalones, separación de regiones por valles, cruce con base de
datos de licencias— se implementó como un servicio independiente en Python
(Flask), separado del backend Node.js.

Esta decisión responde a la naturaleza del procesamiento requerido: el análisis
involucra cálculos numéricos intensivos (FFT, integración de potencia,
suavizado Savitzky-Golay, detección de picos, evaluación de cumplimiento)
para los cuales el ecosistema Python ofrece librerías maduras y optimizadas
(NumPy, SciPy, Pandas). Implementar estas mismas rutinas en Node.js habría
requerido reescribir algoritmos complejos o depender de puentes a librerías
nativas con mantenimiento incierto.

La separación como microservicio también permite escalar el postprocesamiento
de forma independiente, asignar más recursos computacionales al servicio
Python sin afectar la capacidad de respuesta del backend, y mantener
ciclos de desarrollo separados para la lógica de análisis espectral.

\subsection{Separación de canales WebSocket para datos y audio}

El audio demodulado se transmite por un canal WebSocket independiente del
canal de datos de sensores. Esta separación permite tratar cada flujo según
sus requisitos específicos:

\begin{itemize}
    \item El canal de datos transmite eventos discretos (estado, GPS,
    espectro) con frecuencia de segundos y payloads de tamaño moderado.
    \item El canal de audio transmite tramas continuas en formato Opus con
    frecuencia de milisegundos, requiriendo menor latencia y mayor throughput.
    \item La separación evita que la carga de streaming de audio afecte la
    entrega de datos de espectro durante un monitoreo, y viceversa.
\end{itemize}

\section{Problemas presentados y ajustes realizados}

A continuación se documentan los problemas identificados durante el desarrollo
y operación de la plataforma, las evidencias recolectadas, el análisis de
causas, las decisiones de ajuste tomadas y los resultados obtenidos.

\begin{longtable}{p{2.7cm} p{2.7cm} p{2.7cm} p{3cm} p{3cm}}
\caption{Problemas presentados, decisiones y ajustes realizados en la
plataforma web.}\\
\toprule
\textbf{Problema} & \textbf{Evidencia} & \textbf{Causa probable} &
\textbf{Decisión o ajuste} & \textbf{Resultado} \\
\midrule
\endfirsthead
\toprule
\textbf{Problema} & \textbf{Evidencia} & \textbf{Causa probable} &
\textbf{Decisión o ajuste} & \textbf{Resultado} \\
\midrule
\endhead

Límite de capacidad de SQLite causaba bloqueos en producción al alcanzar
~2 GB de datos acumulados &
Logs de error del servidor, mediciones de tamaño de archivo de base de datos &
SQLite no está diseñado para volúmenes de series temporales con múltiples
escrituras concurrentes desde sensores &
Migración completa a PostgreSQL con extensión TimescaleDB. Se reescribió la
capa de conexión, modelos de datos y rutas con consultas directas manteniendo
compatibilidad de formatos &
Eliminación del límite de capacidad, mejora de rendimiento en consultas,
habilitación de compresión y retención automática de datos antiguos \\

\addlinespace

Endpoint de campañas para sensores con nombre incorrecto y resultados
duplicados en asignaciones múltiples &
Logs de consulta de sensores físicos, pruebas de integración &
Falta de deduplicación en consulta SQL de asignación sensor-campaña y
nomenclatura inconsistente del endpoint &
Corrección del nombre del endpoint, incorporación de lógica de deduplicación,
adición de filtro opcional por estado y validación de campos requeridos &
Sensores reciben lista única de campañas asignadas con parámetros completos
para ejecutar mediciones \\

\addlinespace

Timeouts en generación de reportes de cumplimiento para campañas con grandes
volúmenes de datos &
Errores HTTP 504 en frontend, logs de timeout en Nginx &
Timeouts por defecto (30--60 s) insuficientes para el pipeline completo de
consulta, postprocesamiento y armado de respuesta &
Configuración de timeouts extendidos a 3600~s en servidor Node.js y proxy
Nginx (proxy\_read\_timeout, proxy\_connect\_timeout, proxy\_send\_timeout) &
Reportes generados exitosamente incluso para campañas con grandes volúmenes
de datos acumulados \\

\addlinespace

Sensores que dejaban de reportar permanecían con estado activo en la
plataforma &
Reportes de personas usuarias, inspección de estados inconsistentes en base
de datos &
Ausencia de mecanismo automático de verificación de frescura de datos de
sensores &
Implementación de tarea programada cada 30~s que marca como offline sensores
sin reporte reciente y notifica cambios al frontend por WebSocket &
El estado de sensores refleja fielmente su condición operativa real,
eliminando falsos positivos \\

\addlinespace

Carga inicial lenta del frontend por inclusión de Plotly.js en el bundle
principal &
Métricas de Lighthouse, reportes de personas usuarias sobre tiempo de carga &
Plotly.js (~3~MB) se incluía en el bundle principal, bloqueando renderización &
Separación de Plotly.js en chunk independiente mediante code splitting en
Vite; pantallas sin gráficos no esperan su carga &
Mejora significativa en tiempo de carga inicial y puntuación de rendimiento
sin afectar pantallas de monitoreo \\

\addlinespace

Errores de Mixed Content en producción con HTTPS: conexiones WebSocket y API
bloqueadas por navegador &
Errores en consola del navegador, fallos en conexión WebSocket, datos en
tiempo real sin actualizar &
URLs del backend configuradas como absolutas (HTTP) desde página servida por
HTTPS &
Migración a URLs relativas (\texttt{/api}, \texttt{/ws}) en variables de
entorno de compilación, delegando resolución al proxy Nginx &
Eliminación de errores de Mixed Content, funcionamiento correcto en todos los
entornos \\

\addlinespace

Inconsistencia de sesión entre múltiples pestañas del navegador &
Reportes de personas usuarias, pruebas con múltiples pestañas abiertas &
Estado de autenticación y WebSocket gestionado por pestaña sin coordinación &
Revisión del flujo de MSAL para propagar eventos de login/logout entre
pestañas y manejo de renovación de tokens &
Comportamiento consistente de sesión independientemente del número de
pestañas abiertas \\
\bottomrule
\end{longtable}

\section{Validación, pruebas o evidencias}

La plataforma web ha sido validada mediante múltiples estrategias de
verificación a lo largo de su desarrollo y operación:

\textbf{Pruebas funcionales manuales}: cada pantalla del frontend y cada
endpoint de la API han sido verificados manualmente durante el desarrollo,
cubriendo los flujos principales de autenticación, monitoreo en vivo, creación
de campañas, revisión de alertas y generación de reportes.

\textbf{Pruebas de integración con sensores físicos}: se ha verificado la
comunicación extremo a extremo con sensores HackRF reales, incluyendo el envío
de configuraciones, la recepción de datos de espectro y estado, y la
transmisión de audio demodulado AM/FM. El simulador de sensores AM/FM
documentado en el backend permite reproducir escenarios de prueba sin hardware
físico.

\textbf{Documentación interactiva de API}: la especificación OpenAPI publicada
a través de Swagger UI permite verificar la estructura de cada endpoint, los
esquemas de solicitud y respuesta, y ejecutar pruebas directamente desde el
navegador mediante la función ``Try it out''.

\textbf{Validación en producción}: la plataforma se encuentra en operación en
el servidor institucional de la ANE, accesible a través de la red interna. El
uso continuo por parte del personal técnico constituye una validación
operativa de las funcionalidades implementadas.

\textbf{Logs de servidor}: el backend registra cada solicitud HTTP con método,
ruta, marca de tiempo y código de respuesta, permitiendo auditoría de la
actividad y diagnóstico de problemas. Los logs de Nginx complementan esta
trazabilidad con información de tiempos de respuesta y estado de las
conexiones.

\textbf{Mantenimiento automático verificado}: la tarea programada de
validación de estado de sensores se monitorea a través de los eventos
WebSocket que notifican cambios de estado al frontend. Los sensores que dejan
de reportar son marcados como offline de forma automática, y este
comportamiento ha sido confirmado en operación.

\section{Aprendizajes y transferencia de conocimiento}

Los siguientes aprendizajes se derivan de la experiencia de diseño,
implementación y operación de la plataforma web:

\begin{enumerate}
    \item \textbf{Dimensionar la persistencia desde el inicio según la carga
    esperada}: SQLite fue adecuado para desarrollo y prototipado, pero
    mantener una base de datos sin servidor para una carga de producción con
    múltiples escritores concurrentes y grandes volúmenes de series temporales
    generó un cuello de botella que requirió una migración significativa. La
    enseñanza es que la elección de la base de datos debe basarse en el perfil
    de carga previsto —volumen, concurrencia, tipo de consultas— y no solo en
    la conveniencia de desarrollo.

    \item \textbf{La comunicación en tiempo real debe ser push, no pull}: el
    uso de WebSocket para notificar al frontend sobre nuevos datos y cambios
    de estado evitó la necesidad de implementar lógica de polling, reduciendo
    la carga en el servidor y la latencia de visualización. Para sistemas
    donde los datos fluyen de forma continua desde múltiples fuentes, el
    patrón push es la estrategia correcta.

    \item \textbf{Los timeouts deben calibrarse según la operación más
    costosa}: los valores por defecto de los servidores HTTP (30 segundos) no
    contemplan operaciones de análisis que involucran consultas complejas,
    procesamiento numérico y cruce con bases de datos externas. Identificar la
    operación de mayor duración —los reportes de cumplimiento— y configurar
    todos los timeouts de la cadena (servidor, proxy) de forma consistente
    evita fallos intermitentes difíciles de diagnosticar.

    \item \textbf{La separación de canales por tipo de dato simplifica el
    manejo de calidades de servicio}: mantener datos de espectro y audio en
    canales WebSocket independientes permitió tratar cada flujo con sus
    propios requisitos de latencia y frecuencia, sin que la carga de
    streaming de audio afectara la entrega de datos de espectro.

    \item \textbf{Las URLs relativas eliminan problemas de configuración por
    entorno}: migrar de URLs absolutas a relativas para API y WebSocket
    resolvió permanentemente los errores de Mixed Content y permitió que una
    misma compilación funcione en desarrollo, VPN y producción sin
    reconfiguración.

    \item \textbf{La división de código es necesaria en aplicaciones con
    dependencias de visualización pesadas}: las librerías de gráficos
    científicos pueden dominar el tamaño del bundle. Planificar code splitting
    desde el inicio evita resolver problemas de rendimiento cuando la
    aplicación ya está en producción.

    \item \textbf{El mantenimiento automático del estado elimina falsos
    positivos}: sin verificación periódica de la frescura de datos de
    sensores, los equipos desconectados aparecían como operativos. La
    automatización de esta verificación, con notificación por WebSocket,
    eliminó este problema y redujo la carga de monitoreo manual del personal.

    \item \textbf{La documentación interactiva de API acelera la integración}:
    contar con Swagger UI permitió que el equipo de frontend y los
    desarrolladores de los sensores pudieran explorar y probar la API sin
    depender de documentación estática, reduciendo el tiempo de integración.
\end{enumerate}

\section{Limitaciones}

La plataforma web en su estado actual presenta las siguientes limitaciones:

\begin{enumerate}
    \item \textbf{Base de datos sin alta disponibilidad}: PostgreSQL opera en
    instancia única, sin replicación ni mecanismo de failover automático. Una
    falla del servidor de base de datos interrumpiría la totalidad de la
    plataforma.

    \item \textbf{Escalamiento vertical del backend}: el backend está diseñado
    como un solo proceso Node.js. Si bien el pool de 20 conexiones soporta la
    concurrencia actual, la arquitectura no contempla escalamiento horizontal
    a múltiples instancias, lo que requeriría un backend compartido para
    WebSocket (como Redis) y balanceo de carga.

    \item \textbf{Sin balanceo de carga}: el punto único de entrada a través
    de Nginx no está configurado para distribuir tráfico entre múltiples
    réplicas del backend.

    \item \textbf{Dependencia de conectividad de sensores}: el modelo de
    operación asume que los sensores consultan periódicamente su configuración.
    Si un sensor pierde conectividad en el momento de consulta, puede no
    recibir una campaña programada hasta el siguiente ciclo.

    \item \textbf{Postprocesamiento síncrono sin cola}: la generación de
    reportes invoca al servicio Python de forma síncrona. Si el servicio no
    está disponible, el reporte falla sin mecanismo de reintento automático o
    encolamiento.

    \item \textbf{Tablas de series temporales sin optimización completa}:
    solo \texttt{sensor\_status} está configurada como hypertable de
    TimescaleDB. Las tablas \texttt{sensor\_data} y \texttt{sensor\_gps},
    que también almacenan series temporales, pueden presentar degradación al
    crecer en volumen.

    \item \textbf{Tamaño de dependencias de visualización}: Plotly.js (~3 MB)
    sigue siendo una dependencia significativa que afecta la primera visita a
    las pantallas de monitoreo, incluso estando en un chunk separado.

    \item \textbf{Sin pruebas automatizadas}: el desarrollo no incluye pruebas
    unitarias ni de integración automatizadas, dependiendo de validación
    manual para verificar cambios.

    \item \textbf{Sin modo offline}: la aplicación requiere conectividad
    constante con el backend. No hay almacenamiento local de datos para
    consulta sin conexión.

    \item \textbf{Sin diseño responsivo para dispositivos móviles}: la
    interfaz está optimizada para estaciones de escritorio, que es el contexto
    de uso principal. No se ha priorizado la adaptación a pantallas pequeñas.
\end{enumerate}

\section{Recomendaciones específicas}

Con base en las limitaciones identificadas y la experiencia acumulada, se
proponen las siguientes recomendaciones para la evolución de la plataforma:

\begin{enumerate}
    \item \textbf{Alta disponibilidad de base de datos}: configurar
    replicación streaming en PostgreSQL con al menos una réplica en caliente,
    permitiendo continuidad operativa ante falla del servidor principal.

    \item \textbf{Escalamiento horizontal del backend}: evaluar la
    introducción de Redis como backend compartido para WebSocket y como caché
    de datos frecuentes, habilitando múltiples instancias del servidor Node.js
    coordinadas detrás de un balanceador de carga.

    \item \textbf{Optimización de todas las tablas de series temporales}:
    migrar \texttt{sensor\_data} y \texttt{sensor\_gps} a hypertables de
    TimescaleDB, y configurar políticas de compresión y retención automática
    para todas las tablas de series temporales.

    \item \textbf{Resiliencia del postprocesamiento}: implementar una cola de
    trabajos (por ejemplo, Bull con Redis) para las solicitudes de reportes,
    con reintentos automáticos con backoff exponencial cuando el servicio
    Python no esté disponible.

    \item \textbf{Pruebas automatizadas}: introducir pruebas unitarias con
    Vitest/Jest para el frontend y backend, y pruebas de integración para los
    flujos críticos (autenticación, monitoreo en vivo, creación de campañas,
    generación de reportes).

    \item \textbf{Evaluación de alternativas de visualización}: analizar la
    viabilidad de reemplazar Plotly.js por ECharts o una implementación basada
    en Canvas/WebGL para reducir el tamaño del bundle sin perder capacidades
    de visualización.

    \item \textbf{Monitoreo operacional}: instrumentar el backend con métricas
    de Node.js (consumo de memoria, event loop lag, latencia de respuesta,
    conexiones activas) y exportarlas a un sistema de monitoreo institucional.

    \item \textbf{Internacionalización preparatoria}: si bien el contexto
    actual es exclusivamente colombiano, estructurar los textos de la interfaz
    de forma que una futura traducción no requiera reescribir componentes,
    utilizando una capa de localización (i18n).

    \item \textbf{Revisión de seguridad}: auditar los endpoints de
    administración para verificar que todos requieran autenticación y
    autorización por rol, y documentar explícitamente el modelo de seguridad
    (endpoints públicos de ingesta, endpoints autenticados de administración).

    \item \textbf{Documentación operativa}: elaborar guías de operación para
    el personal técnico de la ANE, cubriendo procedimientos de despliegue,
    respaldo, recuperación y diagnóstico de problemas comunes.
\end{enumerate}

\section{Repositorios}

El software de la plataforma web está organizado en dos módulos dentro del
repositorio del sistema. A continuación se presentan las fichas técnicas de
cada componente.

\begin{table}[H]
\centering
\caption{Repositorio asociado al frontend de la plataforma web.}
\label{tab:repo_frontend}
\begin{tabularx}{\textwidth}{>{\bfseries}l X}
\toprule
Nombre del software & ANE Frontend \\
Repositorio & Repositorio pendiente de publicación \\
Rama, versión o commit & main \\
Responsable & Equipo de desarrollo de plataforma \\
Función dentro del sistema & Interfaz web para la operación y administración
del sistema de monitoreo espectral \\
Entradas & Interacciones de la persona usuaria, datos en tiempo real por
WebSocket, respuestas de API REST \\
Salidas & Visualizaciones de espectro y mapas, formularios de configuración,
reportes de cumplimiento, streaming de audio \\
Dependencias & React 19, TypeScript, Vite 5, Tailwind CSS 3, Plotly.js,
Leaflet, MSAL, Axios, Nginx \\
Estado actual & En producción en servidor institucional privado \\
\bottomrule
\end{tabularx}
\end{table}

\begin{table}[H]
\centering
\caption{Repositorio asociado al backend de la plataforma web.}
\label{tab:repo_backend}
\begin{tabularx}{\textwidth}{>{\bfseries}l X}
\toprule
Nombre del software & ANE Backend \\
Repositorio & Repositorio pendiente de publicación \\
Rama, versión o commit & main \\
Responsable & Equipo de desarrollo de plataforma \\
Función dentro del sistema & Centro de coordinación: API REST, WebSocket,
persistencia, gestión de campañas, alertas y reportes \\
Entradas & Solicitudes HTTP del frontend y sensores; datos de espectro, GPS,
estado y audio desde sensores RF \\
Salidas & Respuestas JSON, eventos WebSocket, reportes de cumplimiento
normativo \\
Dependencias & Node.js 18+, Express, TypeScript, PostgreSQL 13+, TimescaleDB
(opcional), ws, Swagger/OpenAPI, Azure AD, servicio Python de
postprocesamiento \\
Estado actual & En producción en servidor institucional privado \\
\bottomrule
\end{tabularx}
\end{table}
